% ===== BOT 06: expand_undersized_gs — giãn giải tích Gaussian dưới cỡ =====

% --- Frame 1: Tổng quan + trạng thái thật ---
\begin{frame}{\texttt{expand\_undersized\_gs}: nhánh "giãn giải tích" --- và sự thật là nó \emph{không chạy}}
\footnotesize
\begin{columns}[T]
\begin{column}{0.46\textwidth}
\begin{itemize}
    \item SADGS có \textbf{ba} thao tác thích nghi theo tần số: \textbf{split} (Gaussian quá to so với texture), \textbf{prune} (đóng góp thấp), và \textbf{expand} (Gaussian \emph{quá nhỏ} so với texture $\Rightarrow$ lãng phí, nhiều điểm thừa).
    \item Hàm \texttt{expand\_undersized\_gs} nằm ở \texttt{scene/gaussian\_model.py:833--864}, hoàn chỉnh và có thể chạy.
    \item \textbf{Nhưng} lời gọi duy nhất trong \texttt{train.py:356--359} \textbf{đã bị comment out}. Mặc định SADGS \textbf{chỉ} SPLIT/PRUNE.
    \item Tham số \texttt{tau\_expand} vẫn được khai báo (\texttt{arguments/\_\_init\_\_.py:135}, giá trị $1.0$) nhưng \textbf{không ảnh hưởng} tới kết quả huấn luyện hiện tại.
\end{itemize}
\end{column}
\begin{column}{0.54\textwidth}
\centering
\includegraphics[width=0.99\linewidth]{sadgsx/06_timeline_commented.png}
\captionof{figure}{\footnotesize Vị trí nhánh bị vô hiệu hoá trong khối densification của \texttt{train.py}.}
\end{column}
\end{columns}
\end{frame}

% --- Frame 2: eta là gì + điều kiện "dưới cỡ" ---
\begin{frame}{Điều kiện chọn Gaussian "dưới cỡ": $\eta < \tau_{expand}$}
\scriptsize
Chỉ báo tần số $\eta$ (chế độ \texttt{wavelength}, \texttt{utils/freq\_utils.py:335--348}) là \textbf{tỉ số không thứ nguyên} giữa kích thước Gaussian chiếu lên ảnh và bước sóng texture cục bộ:
\[
\lambda_{\min}=\frac{1}{\sqrt{\lambda_1(S)}+10^{-5}},\qquad
\eta_{3ch}=\frac{\ell_{\text{truc}}}{\lambda_{\min}},\qquad \ell_{\text{truc}}=\sqrt{u^2+v^2}
\]
với $\lambda_1(S)$ là trị riêng lớn nhất của structure tensor 2D tại điểm chiếu. $\eta>1\Rightarrow$ Gaussian \emph{to hơn} chi tiết (aliasing, cần split); $\eta<1\Rightarrow$ Gaussian \emph{nhỏ hơn} chi tiết (dư điểm, đáng giãn).
\[
\texttt{undersized\_mask}=(\texttt{max\_eta\_3ch}<\tau_{expand})\ \wedge\ (\texttt{max\_eta\_3ch}>0)
\qquad \text{(\texttt{gaussian\_model.py:846})}
\]
\begin{itemize}\itemsep1pt
    \item Mask là \textbf{per-axis} trên tensor $[N,3]$: mỗi trục chính của một Gaussian được xét độc lập.
    \item Điều kiện $>0$ loại các Gaussian chưa từng được quan sát (accumulator còn $0$) --- nếu không chúng sẽ bị giãn vô hạn.
    \item $\tau_{expand}=1.0$ chính là ngưỡng Nyquist, trùng với \texttt{TAU\_HIGH}$=1.0$ dùng để phân loại high-$\eta$ (\texttt{freq\_utils.py:395}).
\end{itemize}
\begin{center}
\includegraphics[width=0.78\linewidth]{sadgsx/06_tau_expand_region.png}\\
\captionof{figure}{\footnotesize Vùng $\eta$ bị tác động ứng với vài giá trị $\tau_{expand}$; ngoài vùng đó hệ số giãn $=1$ (không đổi).}
\end{center}
\end{frame}

% --- Frame 3: công thức cập nhật log-scale ---
\begin{frame}{Công thức cập nhật: cộng thẳng vào \emph{log-scale}}
\scriptsize
Vì $\sigma=\exp(\texttt{\_scaling})$ (\texttt{gaussian\_model.py:39}, \texttt{scaling\_activation = torch.exp}), việc nhân scale trở thành phép \textbf{cộng} trong không gian log:
\[
\boxed{\ \Delta_{\log} = -\tfrac{1}{2}\,\ln\!\big(\max(\eta,10^{-6})\big)\ }
\qquad\Longrightarrow\qquad
\frac{s_{\text{mới}}}{s_{\text{cũ}}}=e^{\Delta_{\log}}=\eta^{-1/2}
\]
\begin{itemize}\itemsep1pt
    \item Hệ số là \textbf{$-0.5$} và \textbf{dấu âm}: \texttt{delta\_log\_scale = -0.5 * torch.log(eta\_vals)} (\texttt{gaussian\_model.py:855}).
    \item Với $\eta<1$ thì $\ln\eta<0\Rightarrow\Delta_{\log}>0\Rightarrow$ \textbf{scale tăng} (giãn). Đúng chiều mong muốn.
    \item \texttt{torch.clamp(...,\,min=1e-6)} (\texttt{:854}) chặn trên biên độ: $\Delta_{\log}\le -0.5\ln 10^{-6}\approx 6.91$, tức hệ số giãn tối đa $\approx\!1000\times$ --- vẫn là một con số rất lớn.
    \item Cập nhật viết vào tensor thưa: \texttt{delta = zeros\_like(\_scaling)}; \texttt{delta[undersized\_mask] = delta\_log\_scale}; \texttt{new\_scaling = \_scaling + delta} (\texttt{:857--860}) --- các trục không thoả mask giữ nguyên tuyệt đối.
\end{itemize}
\begin{center}
\includegraphics[width=0.52\linewidth]{sadgsx/06_scale_ratio_eta.png}\\
\captionof{figure}{\footnotesize Hệ số giãn $s_{\text{mới}}/s_{\text{cũ}}=\eta^{-1/2}$ (log--log), kèm ngưỡng $\tau_{expand}$ và biên clamp $10^{-6}$.}
\end{center}
\end{frame}

% --- Frame 4: vì sao gọi là "giải tích" ---
\begin{frame}{Vì sao gọi là "giải tích" (analytic) --- và vì sao nó \emph{miễn phí} bộ nhớ}
\scriptsize
\begin{columns}[T]
\begin{column}{0.44\textwidth}
\begin{itemize}\itemsep2pt
    \item Không lấy mẫu, không sinh con, không đợi gradient hội tụ: từ $\eta$ đo được, \textbf{nghiệm đóng} $\sigma^\ast=\sigma\,\eta^{-1/2}$ được áp \textbf{một lần}.
    \item Không gọi \texttt{densification\_postfix}, không \texttt{cat} tensor $\Rightarrow$ $N$ \textbf{không đổi}; toàn bộ accumulator (\texttt{max\_eta\_3ch}, \texttt{densify\_count}, \dots) giữ nguyên chỉ số.
    \item Đối lập với \texttt{densify\_and\_split\_structgs} (\texttt{:698}): $k=\lceil\sqrt{\eta}\rceil$ mỗi trục, sinh $N_{\text{con}}=k_xk_yk_z$ Gaussian mới.
    \item Chi phí duy nhất: một lần ghi đè tensor \texttt{\_scaling} $[N,3]$.
\end{itemize}
\end{column}
\begin{column}{0.56\textwidth}
\centering
\includegraphics[width=0.99\linewidth]{sadgsx/06_ellipse_expand.png}\\
\captionof{figure}{\footnotesize Ellipse chiếu trước/sau khi giãn trên nền texture bước sóng $\lambda$: chỉ tham số đổi, số Gaussian không đổi.}
\end{column}
\end{columns}
\end{frame}

% --- Frame 5: bất nhất docstring vs code ---
\begin{frame}{Phát hiện: docstring giả định $\eta\propto(\sigma\omega)^2$, code thực tế đo $\eta$ \emph{tuyến tính}}
\scriptsize
Docstring (\texttt{gaussian\_model.py:837--840}) viết: \emph{"eta = (sigma * omega)$^2$ \dots\ log(sigma\_target) = log(sigma\_old) - 0.5 * log(eta)"}, tức giả định $\eta$ \textbf{bậc hai} theo scale, khi đó một bước đưa $\eta\to 1$ \emph{chính xác}.
Nhưng \texttt{freq\_utils.py:348} tính $\eta=\ell_{\text{truc}}/\lambda_{\min}$ --- \textbf{bậc nhất} theo scale. Thay vào:
\[
\eta_{\text{mới}}=\frac{\ell\cdot \eta^{-1/2}}{\lambda_{\min}}=\eta\cdot\eta^{-1/2}=\sqrt{\eta}\ \ne\ 1
\qquad\Longrightarrow\qquad \eta_k=\eta_0^{(1/2)^k}
\]
\begin{itemize}\itemsep1pt
    \item Hệ quả: một lần gọi \textbf{không} đạt Nyquist như comment "\texttt{This makes eta\_new = 1.0 exactly}" (\texttt{:853}) tuyên bố --- chỉ \textbf{rút căn} khoảng cách log tới $1$.
    \item Hội tụ vẫn rất nhanh (bậc hai trong không gian $\ln\eta$): $\eta_0=10^{-4}$ cần $\approx 4$ lần gọi để lên $>0.5$.
    \item Lưu ý chế độ \texttt{projection} (\texttt{freq\_utils.py:371--373}) cũng lấy \texttt{sqrt} nên cũng tuyến tính theo scale $\Rightarrow$ kết luận không đổi.
\end{itemize}
\begin{center}
\includegraphics[width=0.76\linewidth]{sadgsx/06_eta_convergence.png}\\
\captionof{figure}{\footnotesize Trái: $\eta$ sau 1 lần gọi (giả định docstring vs code thật). Phải: dãy $\eta_k=\eta_0^{(1/2)^k}$.}
\end{center}
\end{frame}

% --- Frame 6: ghi vào _scaling qua optimizer ---
\begin{frame}{Ghi kết quả vào \texttt{\_scaling}: đi \emph{qua} optimizer, không gán trực tiếp}
\footnotesize
\begin{itemize}\itemsep2pt
    \item Toàn bộ phép giãn nằm trong \texttt{with torch.no\_grad():} (\texttt{:851}) --- đây là can thiệp ngoài đồ thị tính đạo hàm.
    \item Không gán \texttt{self.\_scaling.data = \dots}, mà gọi
    \texttt{optimizable\_tensors = self.replace\_tensor\_to\_optimizer(new\_scaling, "scaling")} rồi
    \texttt{self.\_scaling = optimizable\_tensors["scaling"]} (\texttt{:863--864}).
    \item Trong \texttt{replace\_tensor\_to\_optimizer} (\texttt{:485--501}): tìm \texttt{param\_group} tên \texttt{"scaling"}, \textbf{đặt lại} \texttt{exp\_avg} và \texttt{exp\_avg\_sq} (và \texttt{max\_exp\_avg\_sq} nếu có, do \texttt{optimizer\_type="hybrid"} bật \texttt{amsgrad=True} tại \texttt{:323}) về $0$, xoá state cũ, bọc tensor mới thành \texttt{nn.Parameter} và gắn lại state.
    \item \textbf{Điểm đau}: moment Adam bị xoá cho \textbf{toàn bộ} $N$ Gaussian, kể cả các Gaussian không hề bị giãn --- một cú sốc động lượng toàn cục mỗi \texttt{densification\_interval}.
    \item \textbf{Rủi ro kỹ thuật}: \texttt{stored\_state} được dùng ngay mà \emph{không} kiểm tra \texttt{None} (\texttt{:489--490}) $\Rightarrow$ nếu gọi trước bước \texttt{optimizer.step()} đầu tiên sẽ \texttt{AttributeError}.
\end{itemize}
\end{frame}

% --- Frame 7: chi phí bộ nhớ ---
\begin{frame}{Expand vs.\ "tile bằng bản sao cũ": chênh lệch bộ nhớ bậc ba}
\scriptsize
\texttt{densify\_and\_split\_structgs} chỉ kích hoạt khi $\eta>1$ (Gaussian \emph{quá to}) --- nó
\textbf{không hề chạy} cho Gaussian dưới cỡ ($\eta<1$). Vậy nếu tắt \texttt{expand}, cách duy nhất để một
vùng dưới cỡ đạt cùng độ phủ là \textbf{xếp nhiều bản sao giữ nguyên kích thước cũ} lấp đầy chỗ đó --- đây
là một phép \emph{ước lượng bậc độ lớn} minh hoạ, \textbf{không phải} hành vi thật của bất kỳ hàm nào trong code:
\[
\text{expand: } N\mapsto N\ (\Delta=0)
\qquad\text{vs.}\qquad
\text{"tile bằng bản sao cũ": } N\mapsto N\cdot k^3,\quad k=\Big\lceil\frac{1}{\eta}\Big\rceil
\]
(hệ số $1/\eta$, \textbf{không phải} $\eta^{-1/2}$: theo Frame 5, $\eta$ đo được là tỉ số \emph{tuyến tính}
theo scale, nên cần nhân đúng $1/\eta$ để đưa $\eta_{\text{mới}}\to 1$; luỹ thừa $3$ vì đang lấp một \emph{thể tích} 3D theo cả 3 trục.)
\begin{itemize}\itemsep1pt
    \item Mỗi Gaussian chiếm $59$ float32 $=236$\,B (\texttt{xyz}\,3 $+$ \texttt{scaling}\,3 $+$ \texttt{rotation}\,4 $+$ \texttt{opacity}\,1 $+$ \texttt{f\_dc}\,3 $+$ \texttt{f\_rest}\,45), chưa kể trạng thái Adam.
    \item $\eta=0.1\Rightarrow k=10\Rightarrow k^3=1000\times$ số Gaussian để lấp cùng thể tích; $\eta=0.01\Rightarrow k=100\Rightarrow k^3=10^6\times$.
    \item Đây là lý do thiết kế cho \texttt{expand}: một \textbf{van giảm áp VRAM tiềm năng} của SADGS --- nhưng vì nó đang bị tắt (Frame 1), lợi ích này \textbf{không có thật} trong pipeline đang chạy.
\end{itemize}
\begin{center}
\includegraphics[width=0.82\linewidth]{sadgsx/06_memory_expand_vs_clone.png}\\
\captionof{figure}{\footnotesize Ước lượng bậc độ lớn: bộ nhớ tham số để lấp cùng một thể tích bằng bản sao
kích thước cũ ($N\cdot k^3$, $k=\lceil 1/\eta\rceil$) so với expand ($\Delta N=0$). Đây không phải hành vi
của \texttt{densify\_and\_split\_structgs}, vốn không chạy khi $\eta<1$.}
\end{center}
\end{frame}

% --- Frame 8: rủi ro nếu bật lại ---
\begin{frame}{Rủi ro nếu bật lại nhánh expand}
\footnotesize
\begin{enumerate}\itemsep2pt
    \item \textbf{Sai đầu vào}: \texttt{train.py:358} truyền \texttt{avg\_high\_eta\_3ch}, mà tensor này \emph{chỉ} khác $0$ ở các Gaussian có \texttt{eta\_high\_count}$>0$ (\texttt{train.py:349--350}), tức các Gaussian \textbf{đã} vi phạm Nyquist. Kết hợp với điều kiện $>0$, expand sẽ chỉ chạm đúng nhóm đang chờ \textbf{split} --- giãn và chia cùng một lúc, mâu thuẫn trực tiếp.
    \item \textbf{Nổ scale}: clamp $10^{-6}$ cho phép nhân scale tới $\approx\!1000\times$ trong một lần; Gaussian khổng lồ làm hỏng tile-sorting, tăng thời gian raster và nuốt gradient của hàng xóm.
    \item \textbf{Reset Adam toàn cục} cho param group \texttt{scaling} mỗi lần gọi (xem frame trước) $\Rightarrow$ dao động loss quanh mốc densify.
    \item \textbf{Chống lại prune}: Gaussian $\eta$ thấp cũng là ứng viên prune (\texttt{low\_ratio > prune\_ratio\_threshold}$=0.8$, \texttt{train.py:353}); giãn chúng ngay trước \texttt{densify\_and\_prune\_structgs} là làm việc thừa.
    \item \textbf{Crash sớm}: \texttt{stored\_state} \texttt{None} nếu gọi trước bước optimizer đầu tiên.
    \item \textbf{Chưa hiệu chỉnh}: do $\eta$ tuyến tính, một bước chỉ đạt $\sqrt{\eta}$ --- muốn đúng Nyquist phải dùng $\Delta_{\log}=-\ln\eta$ (hệ số $-1.0$), hoặc lặp nhiều chu kỳ.
\end{enumerate}
\vspace{0.1cm}
\textbf{Kết luận trạng thái:} \texttt{expand\_undersized\_gs} là code \emph{có mặt nhưng chết} --- \textbf{KHÔNG chạy mặc định}. Mọi kết quả SADGS báo cáo hiện tại đều \textbf{không} bao gồm cơ chế giãn giải tích; \texttt{tau\_expand}, \texttt{expansion\_speed}, \texttt{adaptive\_clone} (\texttt{arguments/\_\_init\_\_.py:135--137}) là tham số ngủ đông.
\end{frame}
